\subsection*{Beyond dense FP16: what else the four steps express}

The examples above use \texttt{mma.row.row} on dense FP16 tiles. The same four-step pattern extends to workloads the basic form cannot reach.

\textbf{Structured sparsity.} When half the elements of A follow a fixed sparse zero pattern, the backend can skip the zero products --- but it needs a \textbf{metadata operand} \texttt{me} that encodes which elements are nonzero:

\begin{lstlisting}[language=croqtile,caption={A sparse contraction}]
mma.row.row.sp mc, ma, mb, me;
\end{lstlisting}

The \texttt{.sp} suffix adds the metadata operand; everything else is the same fill-load-multiply-store rhythm. Any layout combination works: \texttt{mma.row.col.sp}, etc. You load \texttt{me} from a separate metadata tensor alongside A and B.

For quantized operands with per-tile scales, the contraction also needs the scale values associated with the selected A and B tiles. The \co@.scale@ form expresses this scaling alongside the matrix operation:

\begin{lstlisting}[language=croqtile,caption={A scaled contraction}]
mma.row.row.scale mc, ma, mb, sc_a, sc_b;
\end{lstlisting}

Each result element is scaled by \texttt{sc\_a} and \texttt{sc\_b} after the contraction --- no separate dequant kernel needed. Alternatively, scaling can be a \textbf{standalone statement} when the scale source differs from the standard fused path:

\begin{lstlisting}[language=croqtile,caption={Another contraction layout}]
mma.row.row mc, ma, mb;
mma.scale mc, sc_a, sc_b;
\end{lstlisting}

The standalone \texttt{mma.scale} appears in some MoE and mixed-precision kernels.

\textbf{Swizzled loads and transposed stores.} When a staging buffer uses a swizzle pattern (Section~\ref{sec:3-6}), the MMA load must use the matching swizzle mode: \texttt{mma.\allowbreak{}load.\allowbreak{}swiz\textless{}N\textgreater{}}. The \texttt{\textless{}N\textgreater{}} must agree between \texttt{dma.\allowbreak{}copy.\allowbreak{}swiz\textless{}N\textgreater{}} and \texttt{mma.\allowbreak{}load.\allowbreak{}swiz\textless{}N\textgreater{}} --- a mismatch reads garbage. For output, \texttt{mma.\allowbreak{}store.\allowbreak{}transp mc, dst} writes the accumulator with rows and columns swapped, useful when the next stage expects column-major data.

\textbf{Pipeline synchronization.} In pipelined kernels where a producer and a consumer region overlap (Section~\ref{sec:3-8-5}), \texttt{mma.commit} marks the boundary between "done reading this K-slab's operands" and "safe to reuse the staged buffer." It is mandatory glue in event-driven pipelines.

These extensions all follow the same design principle: the four-step skeleton stays fixed, and the variant suffix communicates a specific contract to the backend. The table below collects every variant for reference.
